% !TeX root = ../main.tex
% ---------------------------------------------------------------
\section{Data}
\label{sec:data}
% ---------------------------------------------------------------

\subsection{OpenAlex corpus and eligible researchers}

The corpus is drawn from the October 2026 snapshot of OpenAlex \citep{priemOpenAlexFullyopenIndex2022}. It contains works published from 2000 to 2025 in journals, conference proceedings and book series, of nine work types, with at most 29 authors, after removal of duplicate titles. The spectral rankings of sources and institutions use the works of 2020--2024 and are computed for each of the 26 OpenAlex fields. Work scores are then attached to every corpus work whose source is retained in at least one field, whatever its publication year, so a paper of 2005 is scored by the present-day influence of its journal and institutions and by the influence of the sources that cite it.

Authors are identified by the OpenAlex author identifier. Following the volume baseline of the Stanford database, which reports percentiles for every author with at least five papers \citep{ioannidisStandardizedCitationMetrics2019a}, we treat as eligible every author with at least five corpus works of type article, review or conference paper. Of the \NOaAuthors{} authors with at least one corpus work, \NOaEligible{} are eligible. For comparison, the field tables of the Stanford database count \NStPopul{} Scopus authors with at least five papers. The difference probably reflects Scopus's longer time span, its broader coverage of document types and large collaborations, and differences in author disambiguation. Researcher scores are computed over the eligible authors' articles, reviews and conference papers only, and a score is ranked only when it rests on at least five works: five works with a publication score, or five cited works for a citer score.

Two features of OpenAlex author records bear on the results. Author disambiguation is algorithmic and imperfect, so some researchers are split across several identifiers and some identifiers merge different people. And the corpus omits papers with 30 or more authors, which removes the large-collaboration papers of particle physics, astronomy and genomics that weigh heavily in the Stanford indicators for those fields.

\subsection{The Stanford database}

We use version 9 of the database, released in August 2026 and built from a Scopus snapshot of 1 August 2026 \citep{ioannidisUpdatedSciencewideAuthor2020}. The career table lists \NStCareer{} researchers on citations received from 1996 to 2025 to papers published since 1960. The single-year table lists \NStSingle{} researchers on citations received in 2025. Each row gives the researcher's name, one institution and country, the number of papers, the six component indicators and $c$ with and without self-citations, the self-citation share, and the researcher's main Science-Metrix field and subfields. Scopus author identifiers are not included.

\subsection{Matching Stanford researchers to OpenAlex}

Without a shared identifier, Stanford researchers must be matched to OpenAlex authors by name. We pool the two tables into \NStPersons{} distinct persons. Candidates are eligible OpenAlex authors who have a name variant, on any of their authorships, with the same surname and first initial, whose affiliations include the Stanford country, and whose $h$-index is within a factor of 2.5 of the Stanford career $h$. For persons who appear only in the single-year table, whose $h$ counts 2025 citations alone, the paper count is compared instead. The five candidates closest in paper count and $h$ are then scored. The author's canonical OpenAlex name must contain the Stanford surname and must agree on given names: full given names must coincide once concatenated (so that Zhong Lin matches Zhonglin but Yanfei does not match Yan), and initials must agree as far as both names give them. Candidates whose names agree only in initials are accepted only when one of their OpenAlex affiliations matches the Stanford institution. Among the remaining candidates, the score rewards an institutional match, agreement of middle initials and closeness in volume, and a person is matched when the best candidate leads the next by a clear margin. An OpenAlex author claimed by two persons is kept for the stronger claim. On this rule \NMatchMatched{} persons are matched, covering \NMatchCareerPct{}\% of the career table and \NMatchSinglePct{}\% of the single-year table.

The match can be checked against quantities that both databases measure. Across matched career-table researchers, the Spearman correlation between Stanford's total citation count and the number of corpus citations to the researcher's OpenAlex works is \NCheckNcRho{}, and between the Stanford and OpenAlex $h$-indices it is \NCheckHRho{}. Both are well short of one, because Scopus and OpenAlex differ in coverage and the corpus starts in 2000, but they are high enough to indicate that the matched records mostly describe the same people. A random sample of 100 matches was also inspected by hand (Section~\ref{sec:results}).
